%! Simplifying Algebraic Expressions and Exponent Rules — Unit 1, Lesson 1.1 — Warm-Up (Answer Key)
\documentclass[10pt]{article}
\usepackage{atda-article}
\usepackage{atda-key}   % pulls in -boxes; do NOT also load -boxes
\newcommand{\TallMath}[1]{$\displaystyle #1\rule[-1.4em]{0pt}{3.2em}$}
\setlength{\workrowsep}{6pt}   % handwriting room; moves blank and key together

\begin{document}

\pageheader{Unit 1, Lesson 1.1}{Warm-Up --- Answer Key}
% NAMESTRIP: no \namedateperiod here — mirrors the blank. Teacher-only prose goes
% in the lesson plan as \begin{teachernote}[Warm-Up], never in a key.

\vspace{0.15in}

\begin{notesbox}{Four Moves Today's Lesson Runs On}
\small
Item 1 is the one to slow down on --- writing the factors out is how today's first rule gets
discovered instead of memorized. Show your work on all four.

\begin{enumerate}[leftmargin=1.6em, itemsep=6pt, topsep=4pt]

\item \textbf{Exponents are repeated multiplication.} Write $x^3 \cdot x^4$ out as a string of
$x$'s, count them, and rewrite the product as a single power of $x$.
\begin{work}
  x^3 \cdot x^4 &= (x\cdot x\cdot x)(x\cdot x\cdot x\cdot x) \\
                &= x^7
\end{work}

\item \textbf{Like terms.} Combine: $5m^2-3m+8-2m^2+7m$
\begin{work}
  5m^2-3m+8-2m^2+7m &= (5m^2-2m^2)+(-3m+7m)+8 \\
                    &= 3m^2+4m+8
\end{work}

\item \textbf{Distributing.} Simplify: $-3(2y-5)+4y$
\begin{work}
  -3(2y-5)+4y &= -6y+15+4y \\
              &= -2y+15
\end{work}

\item \textbf{Spiral --- Lesson 1.0.} A rectangular pen uses $30$ feet of fence on three sides
against a barn wall, so its area in square feet is $A(x)=x(30-3x)$. Expand $A(x)$ and write it
in standard form.
\begin{work}
  A(x) &= x(30-3x) \\
       &= 30x-3x^2 \\
       &= -3x^2+30x
\end{work}

\end{enumerate}
\end{notesbox}

\end{document}
